\subsubsection{Coyote time a jeho implementace}

Jak už bylo zmíněno v koncepci, coyote time dovolí hráči skočit i pár snímků poté, co opustil platformu. V \texttt{PlayerController} se to řeší pomocí proměnné \texttt{lastOnGroundTime}, která se při dopadu na zem nastaví na hodnotu \texttt{coyoteTime} z \texttt{PlayerControllerData}. Následně se v každém snímku v metodě \texttt{Update()} tato proměnná snižuje o \texttt{Time.deltaTime}.

V momentě, kdy hráč stiskne skok, se kontroluje, zda je \texttt{lastOnGroundTime} větší než 0. Pokud ano, skok se provede i když hráč technicky už není na zemi. Toto řešení poskytuje hráči malé okno neperfektnosti, kdy se skok stále zaregistruje, i když už by hráč měl správně spadnout.

Hodnota \texttt{coyoteTime} se typicky nastavuje na 0.1 až 0.2 sekundy, což odpovídá přibližně 6-12 snímkům při 60 snímcích za sekundu. Tato doba je dostatečně krátká, aby neovlivnila běžnou hru, ale dostatečně dlouhá, aby zachytila nepatrně opožděný input hráče.

\begin{lstlisting}[language=csharp,caption={Listing 6: Implementace coyote time}]
// In Update() the timer decreases
private void Update()
{
    lastOnGroundTime -= Time.deltaTime;
    lastPressedJumpTime -= Time.deltaTime;
    IsGrounded();
    
    // Set coyote time when grounded
    if (grounded) 
        lastOnGroundTime = data.coyoteTime;
    
    // Check jump with coyote time
    if (lastPressedJumpTime > 0 && lastOnGroundTime > 0 && !pw.IsWallJumping)
        Jump();
}
\end{lstlisting}
